\section{Markers, binary moments, and the one-triple audit}
\label{sec:markers}
\subsection{Exact natural-\texorpdfstring{$S_3$}{S3} transport}
\begin{proposition}\label{prop:marker-transport}
Let $X\leq T\times S_3^g$ be full on $T$ and every marker. Put
$K_X=X\cap C_3^g$, let $\overline X$ be its image in
$T\times C_2^g$, and let $L=\overline X\cap C_2^g$, of dimension
$\ell$. Then
\[
 \overline X\cong T\times L,\qquad
 X/O^2(X)\cong(T/O^2(T))\times C_2^\ell,
 \qquad d_2(X)=d_2(T)+\ell,\qquad \tau_2(X)=\tau_2(T).
\]
For one marker the numbers of possible $X$ with $\ell=1,0$ are
respectively $1$ and $2^{d_2(T)}-1+|\Epi(T,S_3)|$.
\end{proposition}
\begin{proof}
The graph map $T\to C_2^g/L$ factors through $A_2(T)$ and has a
linear lift, proving the first isomorphism. The kernel $K_X$ has odd
order, so inflation is an isomorphism in binary cohomology. The maximal
binary quotient is therefore that of $T\times L$. K\"unneth and the
identical elementary factor show that the kernel defining $\tau_2$
is unchanged. For one marker Goursat's lemma gives: the full product;
one common-$C_2$ graph for each nonzero binary character of $T$; and
one literal graph for each epimorphism $T\to S_3$. These are disjoint.
\end{proof}

If $p_{T,g,\ell}$ counts the actual marker subgroups with this value of
$\ell$, their exact attachment to a critical product $D$, with $f$
fixed points, has weight
\begin{equation}\label{eq:marker-terminal}
 \frac{1}{6^gg!f!}\sum_\ell p_{T,g,\ell}
             N_D((T/O^2(T))\times C_2^\ell).
\end{equation}
This identity does not license replacing an arbitrary complete exterior
projection by its unmarked count.

\subsection{A character moment on the complete binary error}
Let $\mathcal E_N$ be all fixed-point-free $2$-subgroups of $S_{2N}$
with at least one actual orbit outside Table~\ref{tab:critical}, and put
\[
 E_N=\frac{|\mathcal E_N|}{L_{2N}},\qquad
 pE_N=\frac1{L_{2N}}\sum_{Y\in\mathcal E_N}t(Y),
\]
where $t(Y)$ is the number of its actual two-point orbits.
Set $E_0=pE_0=0$.

\begin{proposition}[Complete binary pair moment]\label{prop:binary-moment}
For $N\geq2$,
\begin{equation}\label{eq:binary-moment}
 pE_N\leq[7+(6N)^{1/4}]E_N+r_NE_{N-1},\qquad
 r_N=\frac{N-1}{4}\frac{c_{N-1}G_{N-1}}{c_NG_N}.
\end{equation}
The same estimate holds for the family of all fixed-point-free binary
subgroups, with its corresponding counts.
\end{proposition}
\begin{proof}
Each pair orbit supplies a nonzero character of $Y$. Let $q(Y)$ be
the number of distinct such characters, $D(Y)$ the number of unordered
pairs of equal characters, and $J_4(Y)$ the number of ordered selections
of four pair orbits with independent characters. Then
\[
 t(Y)-q(Y)\leq D(Y),\qquad
 q(Y)\leq7+J_4(Y)^{1/4}.
\]
For the second inequality, choose one original pair orbit for each
distinct character. The span of at most three independent characters
contains at most seven nonzero characters, so each of four successive
choices has at least $q-7$ possibilities. Thus
$J_4(Y)\geq(q-7)^4$ when $q\geq7$; the case $q<7$ is immediate.

There is an injection from groups $Y$ with four ordered independent
pair orbits to groups $H$ with one selected $E_8$ orbit and an ordering
of its four center-pair orbits, at the same degree. For each ordered
pair partition $\mathcal P$, fix a point chart once and use it to place
an extraspecial group $E(\mathcal P)$ on the eight selected points.
Explicitly, inside $C_2^4\rtimes S_4$, take the even flips semidirect
the regular block $V_4$. Its centre has exactly these four pair orbits.
Fix, once for each ordered partition,
an isomorphism $E(\mathcal P)/Z(E(\mathcal P))\cong C_2^4$.
Replace $Y\leq C_2^4\times J$ by its full preimage in
$E(\mathcal P)\times J$. Independence gives the full $E_8$ projection.
The centre of this local projection, its retained pair order, and the
fixed isomorphism recover $Y$.
The complete exterior projection is unchanged, so its noncritical
orbit is retained. There are at most $\lfloor N/4\rfloor$ possible
selected $E_8$ orbits, each with $4!$ possible orders. Consequently
\[
 \sum_{Y\in\mathcal E_N}J_4(Y)\leq 6N|\mathcal E_N|.
\]
Concavity of $x^{1/4}$ gives
$\sum_Yq(Y)\leq[7+(6N)^{1/4}]|\mathcal E_N|$.

Collapse two pair orbits carrying the same character to one pointed
pair orbit. This operation and its inverse preserve a noncritical
orbit. The exact labelled weight ratio is
\[
 \frac{1/(2^2\,2!)}{1/2}=\frac14.
\]
Thus
\[
 \frac1{(2N)!}\sum_{Y\in\mathcal E_N}D(Y)
 =\frac{1}{4(2N-2)!}\sum_{Y\in\mathcal E_{N-1}}t(Y)
 \leq\frac{N-1}{4(2N-2)!}|\mathcal E_{N-1}|.
\]
Divide by $c_NG_N$ and combine the two bounds.
\end{proof}

For $n=2j+1$, put $p_{2j+1}=c_j+c_{j-1}/6$ and
\begin{equation}\label{eq:uv}
 \alpha_j=\frac{c_j}{p_{2j+1}},\quad
 \beta_j=\frac{c_j}{3p_{2j+1}},\quad
 u_j=\alpha_j+\beta_j[7+(6j)^{1/4}],\quad v_j=\beta_jr_j.
\end{equation}
A single fixed point contributes $\alpha_jE_j$.
Replacing a pointed pair by a natural $S_3$ orbit has weight ratio
$(1/6)/(1/2)=1/3$ and contributes $\beta_jpE_j$.
Thus their combined binary-error contribution is at most
$u_jE_j+v_jE_{j-1}$. Section~\ref{sec:analysis} gives
$c_{j-1}/c_j\sim(96j)^{1/4}$, so $u_j\to1$ and $v_j\to0$.
For the recurrence closure one can avoid this sharper asymptotic input:
the elementary coefficient and Gaussian-ratio bounds give, uniformly,
\[
 0\leq u_j\leq4(j+1),\qquad
 0\leq v_j\leq2\Phi^{-2}(j+1)^3.
\]
Either fixed polynomial is absorbed by half of any positive exponential
rate.  Thus the odd transfer introduces no additional analytic frontier.
In particular this is not a strict contraction on unrestricted complete
counts. It will be applied only after $E_j$ has been estimated.

\subsection{The split one-triple incidence}
For a finite group $K$, a nonzero character $\chi:K\to C_3$ is split
if it has a section. This is equivalent to being nonzero on an element
of order three. The nonsplit characters, including zero, are therefore
the annihilator of the span of the order-three elements in $A_3(K)$.
They form a vector space. Write its dimension as $h_3(K)$ and set
\begin{equation}\label{eq:sigma3}
 \sigma_3(K)=3^{d_3(K)}-3^{h_3(K)}.
\end{equation}
This counts oriented split nonzero characters exactly.

\begin{proposition}[Inverse-complement incidence]\label{prop:c1-incidence}
For any collection $\mathcal P$ of pairs $(K,\chi)$ with $K\leq S_m$
and $\chi$ split and nonzero, define
\[
 q(K,\chi)=|\{x\in K:|x|=3,\ \chi(x)=1\}|.
\]
Then
\begin{equation}\label{eq:inverse-complement}
 |\mathcal P|=\sum_{|x|=3}\ 
 \sum_{\substack{N^x=N,\ N\cap\langle x\rangle=1\\
          (N\langle x\rangle,\chi_{N,x})\in\mathcal P}}
 \frac1{|\{y\in Nx:|y|=3\}|},
\end{equation}
where $\chi_{N,x}$ kills $N$ and sends $x$ to $1\in C_3$.
Furthermore, if $P\lhd K$ is a $2$-group, inflation from $K/P$
preserves $d_3$, $h_3$, and $\sigma_3$.
\end{proposition}
\begin{proof}
For each $(K,\chi)$ insert the identity
$1=\sum_{x:|x|=3,\chi(x)=1}q(K,\chi)^{-1}$.
With $N=\ker\chi$, the pair $(N,x)$ determines both $K$ and $\chi$,
and the denominator is precisely the one displayed in
\eqref{eq:inverse-complement}. Every ternary character kills $P$.
A split character upstairs has an order-three representative whose
image is a representative downstairs. Conversely, choose an order-three
element $y$ downstairs on which the character is nonzero and lift it to
$x\in K$. Then $x^3\in P$, so $(x^3)^{2^a}=1$ for some $a$.
The element $z=x^{2^a}$ satisfies $z^3=1$, and its character value is
$2^a\chi(y)\ne0$ in $\mathbb F_3$. Thus $z$ has order three and
the inflated character splits. This proves splitness in both directions
and hence all three equalities.
\end{proof}

These identities are required tests for applying a counting estimate to
the one-external-triple case, customarily denoted $c=1$. One must bound
the sum of \eqref{eq:sigma3}, not an unmarked subgroup sum, and must
retain the reciprocal in \eqref{eq:inverse-complement}. For example,
on the shared-top group $J=\F_4^r\rtimes C_3$ and target
$\F_4^m\rtimes C_3$, above the fixed top map the exact cocycle and
onto-lift counts are
\[
 4^{m(r+1)},\qquad 4^m\prod_{i=0}^{m-1}(4^r-4^i),
\]
respectively. Maschke's theorem gives $H^1(C_3,\F_4^m)=0$;
$\Hom_{C_3}(\F_4^r,\F_4^m)$ has $4^{mr}$ elements, and the
$4^m$ coboundaries are literal distinct graphs. A second mark on the
same $J$ produces a joint moment on that source. It cannot be replaced
by two independently chosen source groups. These exact tests do not
themselves prove that the high-incidence $c=1$ domain has been exhausted.
